% Source: https://artofproblemsolving.com/community/c3h1582862_17_geometry_questions
% Stable identity: aops-topic-1582862-q05
In one revolution the tyre rolls a distance equal to its circumference:

\begin{asy}
usepackage("amsmath");
size(5cm);
defaultpen(fontsize(10pt));
pair O=(0,0);real r=1.6;
draw(circle(O,r),linewidth(1.2));
draw(-r*dir(20)--r*dir(20),Arrows(TeXHead));label("$26.7$ in",(-0.3,0.3),N);
draw((-2.8,-r)--(3.4,-r));
draw((-r,-r-0.35)--(r,-r-0.35),Arrows(TeXHead));label("one turn: $\pi\cdot26.7$ in",(0,-r-0.35),S,fontsize(8pt));
label("\mbox{}",(min(currentpicture,user=true)+max(currentpicture,user=true))/2,black); // keeps the text after the figure black when LaTeX typesets the labels inline
\end{asy}

\[
\pi\cdot26.7\approx83.8805\text{ inches}.
\]
(The arrow is not to scale.) A mile is $5280\cdot12=63{,}360$ inches, so $70{,}000$ miles are $4{,}435{,}200{,}000$ inches, and the number of revolutions is
\[
\dfrac{4{,}435{,}200{,}000}{83.8805}\approx52{,}875{,}251 .
\]
To the nearest hundred, about $\boxed{52{,}875{,}300}$ revolutions (roughly $53$ million).
